\documentclass[11pt,a4paper]{article}
\usepackage[T1]{fontenc}
\usepackage{lmodern,microtype,amsmath,amssymb,amsthm,booktabs,tabularx,array,xcolor}
\usepackage[margin=25mm,top=24mm,bottom=25mm]{geometry}
\usepackage{fancyhdr}
\usepackage[colorlinks=true,linkcolor=blue!35!black,urlcolor=blue!45!black,citecolor=blue!40!black]{hyperref}
\usepackage{bookmark}
\definecolor{ink}{HTML}{16333D}
\definecolor{accent}{HTML}{1A665B}
\usepackage{titlesec}
\titleformat{\section}{\Large\bfseries\color{ink}}{\thesection}{0.65em}{}
\titleformat{\subsection}{\large\bfseries\color{ink}}{\thesubsection}{0.55em}{}
\newtheorem{theorem}{Theorem}[section]
\newtheorem{lemma}[theorem]{Lemma}
\newtheorem{proposition}[theorem]{Proposition}
\newtheorem{corollary}[theorem]{Corollary}
\newcommand{\CT}{\operatorname{CT}}
\newcommand{\Z}{\mathbb Z}
\pagestyle{fancy}\fancyhf{}\renewcommand{\headrulewidth}{0.3pt}
\fancyhead[L]{\footnotesize\color{ink} Differences of Motzkin numbers}
\fancyhead[R]{\footnotesize\color{ink} Research manuscript}
\fancyfoot[C]{\footnotesize\thepage}
\setlength{\headheight}{15pt}
\setlength{\emergencystretch}{2em}
\raggedbottom
\hypersetup{pdftitle={Differences of Motzkin numbers at dyadic indices and at increasing orders},pdfsubject={Motzkin numbers, dyadic valuations and the sign of high logarithmic differences}}
\begin{document}
\begin{center}
{\LARGE\bfseries\color{ink} Differences of Motzkin numbers at dyadic indices and at increasing orders\par}
\vspace{10pt}
{\normalsize Research manuscript \enspace$\cdot$\enspace 15 November 2026\par}
\end{center}
\begin{abstract}
We prove Conjectures 5.2 and 5.3 of Lengyel (2013) on the binary orders
of differences of Motzkin numbers at shifted dyadic indices. For the paired
shifts $2^q$ and $2^q+1$ we determine their common order and
the sharp onset of the formula. Conjecture 5.1 as stated is false. At shift
$195$ its proposed upper bound fails for every odd $n\ge13$.
A generating function for derivatives of the shift recurrence controls the
loss of powers of two and identifies cancellation at the first surviving digit.
For the ordinary
real logarithm we determine the eventual sign of every order of forward
difference at each sufficiently large index. At even indices all signed
differences of order at least two are positive. At odd indices their first
failure occurs at $n/2-(\log n)/(2\log4)+O(1)$. The proof compares two
endpoint contributions and bounds the inverse Laplace transform of the logarithm.
The same endpoint estimate gives simultaneous positivity up to order $c n$ for
all five sequences in Chen, Guo and Wang's almost-infinite log-monotonicity
conjecture. We distinguish the fixed-order consequence accessible from prior
criteria from the growing-order statements.
\end{abstract}
\section{The two differences}
The Motzkin number $M_n$ counts paths of length $n$ from height zero back to
zero, using increments $-1,0,1$ while staying nonnegative. In particular
$M_0=M_1=1$. Choosing the nonhorizontal steps gives
$M_n=\sum_{j=0}^{\lfloor n/2\rfloor}\binom{n}{2j}C_j$, where
$C_j=(j+1)^{-1}\binom{2j}{j}$. We study the same integer sequence under
two operations. For a dyadic index $N=2^n$, the arithmetic question asks
which binary digit first distinguishes $M_{2N+b}$ from $M_{N+b}$.
For ordinary consecutive indices, the analytic question asks for the sign of
\[
 S_{n,r}(a)=(-1)^r\Delta^r\log a_n,\qquad \Delta f_n=f_{n+1}-f_n.
\]
Here $\log$ is the natural logarithm, and the second expression is used on a
positive tail of a sequence.
 
Lengyel stated the complete three-shift valuation pattern as Conjecture~5.3
\cite[pp.~132--135]{lengyel13}. His Theorems~2.1 and 5.6 already cover the
unshifted case and the even-index case for shift one. Section~\ref{sec:lengyel} determines the remaining residues. Conjecture~5.2
asks whether the two shifts $2^q$ and $2^q+1$ eventually give the same
valuation for every $q\ge1$. We introduce an auxiliary coefficient problem.
If $A_j(x),B_j(x)$ express $M_{x+j}$ in terms of $M_x,M_{x-1}$, what are the
residues of $2^q A'_{2^q}(0)$ and $2^q B'_{2^q}(0)$ modulo eight?
Section~\ref{sec:paired} solves this problem for $q\ge4$, controls the higher
Taylor coefficients and recovers the exact orders in the original conjecture.
Section~\ref{sec:counter} uses the same shift expansion at $j=195$ to refute
the proposed universal upper correction of one in Conjecture~5.1.

Chen, Guo and Wang introduced $R(a)_n=a_{n+1}/a_n$. Their unnumbered
concluding conjecture, on five explicitly named sequences, requires
log-convexity of every even iterate and log-concavity of every odd iterate
on a tail that may depend on the iterate \cite[pp.~24--25 of the preprint]{cgw}.
Since $\log R^j(a)_n=\Delta^j\log a_n$, the full condition is equivalent to
$S_{n,j+2}(a)\ge0$ eventually for every fixed $j\ge0$. Any finite set of
orders admits a common tail by taking the maximum of their thresholds.
Section~\ref{sec:linear} proves a common tail for a range of orders growing
linearly with the index. Section~\ref{sec:motzkin} locates the first
exception in an entire Motzkin row.

\section{Exact dyadic differences at the first three shifts}
\label{sec:lengyel}
For the same Motzkin numbers, put $N=2^n$ and
$D_b(n)=M_{2N+b}-M_{N+b}$. Write $v_2$ for the additive
$2$-adic valuation, normalized by $v_2(2)=1$. Lengyel
posed the following three-shift assertion in Conjecture~5.3 of
\cite{lengyel13}. The unshifted input was proved by Lengyel himself
\cite[Theorem~5.6]{lengyel13}, \cite{lengyel14}.

\begin{theorem}[Lengyel's Conjecture 5.3]
For $b=0,1$ and $n\ge2$, the valuation $v_2(D_b(n))$ is $n-1$ for
even $n$ and $n$ for odd $n$. For $b=2$ and $n\ge3$, it is $n$ for
even $n$ and $n-2$ for odd $n$.
\end{theorem}

\begin{proof}
Set $U_x=M_x$ and $Y_x=M_{x-1}$. The binary-digit formulas of
Krattenthaler and M\"uller \cite[Theorem~11]{km18} imply, for $n\ge4$,
\[
 U_N\equiv3\pmod8,\qquad
 Y_N\equiv\begin{cases}4& (n\text{ even}),\\7&(n\text{ odd})\end{cases}
 \pmod8 .
\]
Indeed, $N$ has one digit $1$ and no adjacent pair of ones. For
$N-1$, the first zero bit in position at least four is at position
$n$, bit $n+1$ is zero, the digit sum is $n$, and the number of
adjacent pairs of ones is $n-1$. The two entries of their theorem
for an index congruent to $15\bmod16$ reduce to $4$ for even $n$
and $2n^2+4(n-1)+5\equiv7\pmod8$ for odd $n$.

The case $b=0$ is Lengyel's theorem cited above. For $b=1$, use the
Motzkin recurrence $(x+2)M_x=(2x+1)M_{x-1}+3(x-1)M_{x-2}$ to write
\[
 M_{x+1}=U_x+x\frac{U_x+3Y_x}{x+3}.
\]
Thus $D_1-D_0$ is divisible by $N$ in $\mathbb Z_2$, proving the even
case. For odd $n\ge5$ the quotient $(D_1-D_0)/N$ reduces modulo $2$
to $U_N+Y_N=0\pmod2$. Since $D_0/N$ is odd, so is $D_1/N$.

For $b=2$, two uses of the same recurrence give
\[
 M_{x+2}=A(x)U_x+B(x)Y_x,\qquad
 A(x)=\frac{7x^2+28x+24}{(x+3)(x+4)},\quad
 B(x)=\frac{3x(2x+5)}{(x+3)(x+4)} .
\]
The exact errors in the first-order expansions are
\[
 A(x)-2-\frac{7x}{6}=-\frac{x^2(7x+19)}{6(x+3)(x+4)},
 \qquad B(x)-\frac{5x}{4}=-\frac{x^2(5x+11)}{4(x+3)(x+4)}.
\]
At $x=N,2N$, the respective errors have valuations at least
$2n-3,2n-4$. For $n\ge5$, subtracting and multiplying by $4/N$ gives
a congruence in $\mathbb Z_2$,
\[
 \frac{4D_2}{N}\equiv \frac{8D_0}{N}
 +\frac{14}{3}(2U_{2N}-U_N)+5(2Y_{2N}-Y_N)\pmod8 .
\]
The middle term is $6\bmod8$ for both parities. For even $n$, the
first term equals $4\bmod8$ by the known valuation of $D_0$, and the
last is $5(14-4)\equiv2\pmod8$. Hence $4D_2/N\equiv4\pmod8$.
For odd $n$, the first term vanishes modulo $8$ and the last equals
$5(8-7)\equiv5\pmod8$. Hence $4D_2/N\equiv3\pmod8$.
These give the two asserted valuations. For $n=2,3,4$ the same
integer recurrence starting with $M_0=M_1=1$ gives
\[
\begin{array}{c|rrr}
n&v_2(D_0)&v_2(D_1)&v_2(D_2)\\\hline
2&1&1&0\\3&3&3&1\\4&3&3&4
\end{array}
\]
and completes the proof.
\end{proof}

\noindent
The congruence for $4D_2/N$ also records its leading residue in the
odd case and identifies the cancellation responsible for the two
extra factors of $2$ in the even case. Neither the fixed-modulus
digit formula nor the unshifted valuation alone determines these
shifted orders.
\section{Adjacent shifts by powers of two}\label{sec:paired}
This section establishes Lengyel's Conjecture~5.2 with explicit eventual orders. Put $N=2^n$, $D_j=M_{2N+j}-M_{N+j}$ and $v=v_2$. We use $v(D_0/N)=-1$ for even $n$ and $0$ for odd $n$, and the residues, for sufficiently large $n$,
\[
 U=2M_{2N}-M_N\equiv3\pmod8,\qquad
 W=2M_{2N-1}-M_{N-1}\equiv
 \begin{cases}2&n\text{ even},\\1&n\text{ odd}\end{cases}\pmod8 .
\]
\begin{theorem}[Lengyel's Conjecture 5.2, with the common order]
For every $q\ge1$, the valuations of $D_{2^q}$ and $D_{2^q+1}$ coincide for all sufficiently large $n$. For $q=1$ the common value is $n$ on even $n$ and $n-2$ on odd $n$. For $q=2$ it is $n-1$ on both parities. For every fixed $q\ge3$, the common value is given by
\[
v(D_{2^q})=v(D_{2^q+1})=n-q+
\begin{cases}2,&q\equiv n\pmod 2,\\1,&q\not\equiv n\pmod2 .\end{cases}
\]
For the common-value formulas the least starting index is $3$ for
$q=1,2$, and $q+1$ for every $q\ge3$.

\end{theorem}
\begin{proof}
First assume $q\ge4$. Write $M_{N+j}=A_j(N)M_N+B_j(N)M_{N-1}$, where $A_{-1}=0$, $B_{-1}=1$, $A_0=1$, $B_0=0$. The functions $A_j(x),B_j(x)$ satisfy the same rational recurrence
\[
 (x+j+2)F_j=(2x+2j+1)F_{j-1}+3(x+j-1)F_{j-2}\quad(j\ge1).
\]
For fixed $j$ they are rational and regular at the origin, with $A_j(0)=M_j$ and $B_j(0)=0$ for $j\ge0$. Put $a_j=A'_j(0)$, $b_j=B'_j(0)$, $a(z)=\sum_{j\ge0}a_jz^j$ and $b(z)=\sum_{j\ge0}b_jz^j$. With $O_j(2^n)$ denoting an element of $2^{n-C_j}\mathbb Z_2$ for a constant depending only on $j$, Taylor expansion gives
\[
 D_j/N=M_jD_0/N+a_j U+b_jW+O_j(2^n).
\]

Let $s=(1-2z-3z^2)^{1/2}$ and $M(z)=(1-z-s)/(2z^2)$. Differentiating the recurrence shows that for either $f=a$ or $f=b$,
\[
 zs^2 f'+(2-3z-3z^2)f=h,\qquad
 h_A=1-s^2M(z),\quad h_B=3z.
\]
The term $3z$ comes from $B_{-1}(0)=1$. As $(z^2/s)'/(z^2/s)=(2-3z-3z^2)/(zs^2)$, this gives
\[
 f(z)=\frac{s(z)}{z^2}\int_0^z\frac{w h(w)}{s(w)^3}\,dw.
\]
The formal series $s$, $s^{-3}$ and $h$ are integral since $s=1-z-2z^2M(z)$. Thus the denominator of $a_j$ or $b_j$ divides $\mathrm{lcm}(1,\ldots,j+2)$. For $m=2^q$, $q\ge4$, write $J=h/s^3$. Extracting the coefficient of $z^m$ and multiplying by $m$ gives
\[
 m[z^m]f=\sum_{\ell=2}^{m+2}\frac m\ell [z^{m+2-\ell}]s\,[z^{\ell-2}]J
 \equiv -2[z^{m-2}]J-4M_{m/2}[z^{m/2-2}]J\pmod8 .
\]
Indeed, terms with $v(\ell)\le q-3$ vanish, while those with $v(\ell)=q-2$ have $\ell\le m$ and are killed by the even coefficient of $s$ of index at least $2$. The only remaining denominators are $m,m/2$. Here $M_{m/2}$ is odd for $q\ge4$, by the dyadic Motzkin residues.

Put $H=s^{-3}$ and $K=M/s$. Thus $J_A=H-K$ and $J_B=3zH$. Expanding about $1-z$ yields
\[
 H=\sum_{j\ge0}(2j+1)\binom{2j}{j}\frac{z^{2j}}{(1-z)^{2j+3}},
 \qquad
 K=\frac1{2z^2}\sum_{j\ge1}\binom{2j}{j}\frac{z^{2j}}{(1-z)^{2j}}.
\]
The latter follows from $K=((1-z)/s-1)/(2z^2)$. Since $v\binom{2j}{j}$ is the number of ones in the binary expansion of $j$, only $j=0$ and powers of two remain in $H$ modulo $4$. For $m=2^q$, $q\ge3$, this gives
\[
 [z^{m-2}]H\equiv0,\qquad [z^{m-3}]H\equiv 1+2(q-1)\pmod4.
\]
For the first equality the leading term is $\binom{m}{2}$, divisible by $4$, and all power-of-two indices $j$ involve $\binom{m}{2j+2}$ with $0<2j+2<m$, hence even. For the second, $\binom{m-1}{2}\equiv1\pmod4$ and each of the $q-1$ powers $j\le(m-3)/2$ contributes $2\binom{m-1}{2j+2}\equiv2\pmod4$.

The second expansion gives
\[
 [z^{m-2}]K=\sum_{j=1}^{m/2}\frac12\binom{2j}{j}\binom{m-1}{2j-1}
 \equiv q+(q-1)(q-2)\pmod4.
\]
For even $k\le m$, $\binom{m-1}{k-1}\equiv-1\pmod4$ if $k\le m/2$, and is $1\pmod4$ otherwise. This follows from the product $(-1)^{k-1}\prod_{1\le i<k}(1-m/i)$, in which only $i=m/2$ can affect the residue. For $j=1$, $\binom{2j}{j}/2=1$. At $j=2^u$, $u\ge1$, its residue is $3\pmod4$: the base $j=2$ is direct and the ratio to the value at $2^{u-1}$ is $\prod_{1\le i\le j,\ i\text{ odd}}(1+j/i)\equiv1\pmod4$ for $u\ge2$. The $q$ power-of-two indices thus contribute $-1+(q-2)+3=q\pmod4$. Each index with two ones in its binary expansion contributes $2$, and there are $\binom{q-1}{2}$ such indices below $m/2$. All other terms vanish.

Consequently $[z^{m-2}]J_A$ is $2$ for even $q$ and $3$ for odd $q$ modulo $4$. The parity at the half-size is obtained by replacing $q$ by $q-1$. The corresponding coefficient of $J_B$ is $3(1+2(q-1))$ modulo $4$. Substitution into the two-term extraction formula gives
\[
  (m a_m,m b_m)\equiv
 \begin{cases}
 (0,2),&q\ {\rm even},\\
 (2,6),&q\ {\rm odd}
 \end{cases}\pmod8.
\]
Since $mM_mD_0/N$ vanishes modulo $8$ for $q\ge4$, the expansion of $mD_m/N$ now gives residues $4,2$ for $q$ even and $n$ even, odd, respectively. For $q$ odd it gives $2,4$, respectively. These residues give the stated orders for $D_m$.

Finally, differentiating the recurrence at $j=m+1$ gives
\[
 (m+3)(a_{m+1}-a_m)=m a_m+3m a_{m-1}+2M_m+3M_{m-1}-M_{m+1},
\]
and the analogous formula for $b_j$ omits the three Motzkin terms. Since all derivative denominators up to $m$ have $2$-part at most $m$, both differences are $2$-adic integers as $m+3$ is odd. The recurrence at $x=0$ gives $(m+3)(M_{m+1}-M_m)=m(M_m+3M_{m-1})$, so $v(M_{m+1}-M_m)\ge q$. Hence $(D_{m+1}-D_m)/N$ is $2$-adically integral for all sufficiently large $n$. Since $v(D_m/N)<0$, the two valuations agree. The same integrality argument holds for $q=2,3$ and can be combined with the corresponding small-shift computations.
\paragraph{The three small exponents.}
Here is an independent finite calculation for $q=1,2,3$ in Conjecture~5.2.
It uses a first-order expansion with a fixed shift, so all error estimates
apply to every sufficiently large $n$. In the notation $N=2^n$ and
$D_b=M_{2N+b}-M_{N+b}$, let $d=D_0/N$, $u=2M_{2N}-M_N$ and
$y=2M_{2N-1}-M_{N-1}$. By Lengyel's unshifted congruence
\cite[Theorem~5.6]{lengyel13}, for even $n\ge4$ one has
$d-3/2\in2\mathbb Z_2$. For odd $n\ge3$ one has
$d-1\in2\mathbb Z_2$. The digit evaluations above give $u\equiv3\pmod8$ and
$y\equiv 2$ or $1\pmod8$ according as $n$ is even or odd.

Set $(s,m)=(4,8)$ for $b=2,3,4,5$ and $(s,m)=(8,16)$ for $b=8,9$.
The entries in the last four columns of the table are residues modulo
$m$, obtained by evaluating and differentiating the displayed recurrence.
\[
\begin{array}{c|r|rr|rr}
 b & M_b & s A'_b(0)& s B'_b(0)& sD_b/N\ (n\ \mathrm{even})&
 sD_b/N\ (n\ \mathrm{odd})\\\hline
 2&2&2&5&4&3\\
 3&4&2&7&4&5\\
 4&9&0&6&2&2\\
 5&21&4&2&6&2\\
 8&323&10&6&14&12\\
 9&835&2&14&6&12
\end{array}
\]
For the last two rows the derivative coefficients are even, so the known
residues of $u,y$ modulo $8$ suffice modulo $16$. Each pair of successive
rows has matching valuations for both parities of $n$. This proves the
three exceptional cases, with the exact eventually attained common
valuations $n,n-1,n-2$ for even $n$ in order $q=1,2,3$, respectively,
and $n-2,n-1,n-1$ for odd $n$.
\end{proof}
\paragraph{Effective and sharp onset.}
For $q\ge4$, the paired valuation formula just proved holds for every $n\ge q+1$, and it fails for the shift $m=2^q$ at $n=q$. To bound the Taylor error uniformly, write $F(x,z)=\sum_{j\ge0}F_j(x)z^j$ for either coefficient family and write $f_k=[x^k]F(x,z)$. Multiplication of the shift recurrence by $z^j$ and summation give
\[
 L F=2F_0+x\bigl(-s^2 F+F_0+3zF_{-1}\bigr),\qquad
 L=zs^2\partial_z+(2-3z-3z^2).
\]
Thus $f_1$ is the integral already evaluated and, for $k\ge2$,
\[
 f_k=-\frac{s}{z^2}\int_0^z\frac{w f_{k-1}(w)}{s(w)}\,dw .
\]
On truncation through degree $m+1$, one application of this formula costs at most $q$ powers of two, since the largest possible denominator is $m+3$ and the remaining series are integral. The same bound applies to $f_1$. Hence $v([z^j] f_k)\ge-kq$ for $j\le m+1$. Moreover, in the last integration the only denominator with valuation $q$ is $m$, and for $j=m,m+1$ its coefficient is multiplied by $[z^{j+2-m}]s=-2$. All other denominators have valuation at most $q-1$. This saves one power at the output and gives $v([z^j]f_k)\ge1-kq$ for $j=m,m+1$ and $k\ge1$.

For $n>q$ all Taylor series used here converge at $x=N,2N$ and agree with the rational functions, for instance by their series identities and the absence of zeros in the denominators there. In $mD_j/N$ the total error from $k\ge2$ has valuation at least $n-q+1$, uniformly for $j=m,m+1$. The two linear approximations for $mD_m/N$ and $mD_{m+1}/N$ are congruent modulo $8$ by the derivative-integrality and ordinary recurrence established above. For $n\ge q+2$ the error vanishes modulo $8$. For $n=q+1$ the leading residue is $2\bmod4$, and the error vanishes modulo $4$, proving the asserted uniform onset. At $n=q$, the residue-class-zero case of the binary-digit formula gives $M_{3m}\equiv5$ and $M_{2m}\equiv3\pmod8$, so $v(D_m)=1$ whereas the eventual formula gives $2$.

For the three small exponents the same integral bound gives
$v_2([z^j]f_k)\ge-k L_j$, $L_j=\lfloor\log_2(j+2)\rfloor$.
In the small-exponent table, the error in $sD_j/N$ therefore has
valuation at least $v_2(s)+n-2L_j$ once $n>L_j$. For $q=1,2$ this
is at least three for $n\ge5$; for $q=3$ it is at least four for
$n\ge7$. These thresholds also suffice for the unshifted residue inputs.
For the intervening indices the following entries are the pairs of
valuations at the two shifts, computed by the integer Motzkin recurrence:
\[
\begin{array}{c|ccccc}
q\backslash n &2&3&4&5&6\\\hline
1 &(0,0)&(1,1)&(4,4)&(3,3)&(6,6)\\
2 &(2,3)&(2,2)&(3,3)&(4,4)&(5,5)\\
3 &(2,6)&(1,1)&(2,2)&(4,4)&(4,4)
\end{array}
\]
These entries give the claimed least starting indices for the formulas.
They are thresholds for the common values. Equality of the two
valuations can start earlier.

\paragraph{Lower bounds for arbitrary odd multipliers.}
Let $a$ be any positive odd integer, $b\ge0$, $x=a2^n$, and
$L_b=\lfloor\log_2(b+2)\rfloor$. For every $n\ge2$,
\[
 v_2\bigl(M_{2x+b}-M_{x+b}\bigr)\ge n-L_b.
\]
When $n\le L_b$ this follows by integrality. For $n>L_b$, on truncation
through degree $b$ each integral in the Taylor-coefficient construction
costs at most $L_b$ powers of two because the largest denominator is $b+2$.
Thus $v_2([z^b]f_k)\ge-kL_b$. It follows that each shift derivative
contribution of order $k\ge1$ in $(M_{2x+b}-M_{x+b})/x$ has valuation
at least $(k-1)n-kL_b\ge-L_b$, and the tail tends to zero.
The zeroth contribution is $M_b(M_{2x}-M_x)/x$, of valuation at least
$-1$ by Lengyel's Theorem~2.1 in \cite{lengyel13}, which applies for
all positive odd $a$ and $n\ge2$. This proves the inequality and in
particular the logarithmic lower-bound clause of Conjecture~5.1
for $a=1$, including its bound $n-q$ at $b=2^q-1$.


\section{A counterexample to the upper bound in Lengyel's Conjecture 5.1}\label{sec:counter}

The additional clause of Lengyel's Conjecture~5.1 asserts,
in the present notation, that for each offset $b$ and all sufficiently
large $n$, $v_2(D_b(n))\le n+1$ when $a=1$. One offset violates
this bound at infinitely many indices.

\begin{theorem}
For $b=195$ and every odd $n\ge13$, the integer
$M_{2^{n+1}+195}-M_{2^n+195}$ is divisible by $2^{n+2}$.
\end{theorem}

\begin{proof}
We first determine the fixed residues needed for the unshifted Motzkin
values. We use the Cartier-operator approach to fixed prime-power congruences
of Rowland and Yassawi and its constant-term formulation
\cite{ry,straub}. With $D=1-z(t^{-1}+1+t)$, let $Q=D^{32}$ and
$P=(1-t^2)D^{31}$ over $(\mathbb Z/64\mathbb Z)[z,t,t^{-1}]$.
Then
\[
 M_j=[z^jt^0](1-t^2)/D=[z^jt^0]P/Q.
\]
Indeed, expanding $(1+t+t^{-1})^j$ and taking the indicated
coefficient gives $\sum_k\binom j{2k}\bigl(\binom{2k}k-
\binom{2k}{k-1}\bigr)$, the Catalan sum for $M_j$.
Let $\Lambda_{d,0}$ keep monomials with $z$-exponent congruent to $d$
modulo two and $t$-exponent even, then subtract $d$ from the former and divide both
exponents by two. For $d=0,1$ define $S_d(A)=\Lambda_{d,0}(AQ)$.
Squaring $D$ modulo $2$ gives $D(z,t)^2\equiv D(z^2,t^2)$.
Each repeated squaring gains one power of $2$, hence
$Q(z,t)^2\equiv Q(z^2,t^2)\pmod{64}$. It follows that
$\Lambda_{d,0}(A/Q)=S_d(A)/Q$. The $z$-digits are therefore read from
least significant to most significant. At the final state $A$, the
output is $[z^0t^0]A$ since $Q(0,t)=1$.

Direct multiplication in the indicated finite coefficient ring gives
\[
 S_0^8P=S_0^7P,\qquad S_1^8P=S_1^6P .
\]
The constant coefficients of $S_1 S_0^7 P$, $S_1^6P$ and $S_1^7 P$
are $19,12,7$, respectively. The finite certificate below verifies
the polynomial identities and coefficients. For $n\ge7$ it follows
that $M_{2^n}\equiv19\pmod{64}$ and
$M_{2^n-1}\equiv12,7\pmod{64}$ for even and odd $n$, respectively.

Now use the rational shift functions $A_b,B_b$ defined by the
Motzkin recurrence. To verify the three required residues, the same
certificate evaluates $M_b$ and the derivatives $A_b'(0),B_b'(0)$
at the single fixed index $b=195$ using exact rational numbers. It gives
\[
 M_b\equiv52,\quad 16A'_b(0)\equiv7,\quad
 16B'_b(0)\equiv11\pmod{64}.
\]
Write $N=2^n$. The known unshifted valuation gives $d=D_0/N$ odd
for odd $n$. The just-proved modulus-$64$ values give
$2M_{2N}-M_N\equiv19$ and $2M_{2N-1}-M_{N-1}\equiv17$.
Put $L_{195}=7$. In the last integration step for the coefficient
of degree $195$, the sole denominator of valuation $7$ is $128$ and
its coefficient is multiplied by $[z^{197-128}]s=-2M_{67}$.
Every other denominator costs at most six powers of two. Consequently
$v_2([z^{195}]f_k)\ge1-7k$ for all $k\ge1$. For $n>7$, the tail in
$16D_{195}(n)/N$ starting with $k=2$ has valuation at least
\[
 \min_{k\ge2}\bigl(4+(k-1)n+1-7k\bigr)=n-9 .
\]

Thus, for odd $n\ge15$,
\[
 16D_{195}(n)/N\equiv 16\cdot52 d+7\cdot19+11\cdot17
 \equiv0\pmod{64}.
\]
At $n=13$, exact recurrence gives $v_2(D_{195}(13))=16$.
At $n=11$ it gives $v_2(D_{195}(11))=12$, so $13$ is the first odd
index from which every odd index violates the bound. The contradiction concerns the
stated universal upper correction $1$. The general two-sided estimate
with unspecified, offset-dependent constants is a separate assertion.
\end{proof}

\paragraph{Finite coefficient certificate.}
The following standalone Python code checks the finite polynomial
identities and rational coefficients. The two polynomial cycles in the
proof extend these checks to all binary lengths in the stated ranges.
{\footnotesize\begin{verbatim}
from collections import defaultdict
from fractions import Fraction as F
def mul(a,b):
    c = defaultdict(int)
    for (i,j),u in a.items():
        for (k,h),v in b.items():
            c[i+k,j+h] += u*v
    return {k:v%64 for k,v in c.items() if v%64}
D = {(0,0):1,(1,-1):-1,(1,0):-1,(1,1):-1}
T = {(0,0):1}
for j in range(31): T = mul(T,D)
Q = mul(T,D)
P = mul({(0,0):1,(0,2):-1},T)
def step(A,d):
    return {((i-d)//2,j//2):v for (i,j),v in mul(A,Q).items()
            if i%2 == d and j%2 == 0}
def orbit(d):
    A = [P]
    for j in range(8): A.append(step(A[-1],d))
    return A
A,B = orbit(0),orbit(1)
assert A[8] == A[7] and B[8] == B[6]
assert [step(A[7],1)[0,0],B[6][0,0],B[7][0,0]] == [19,12,7]
m=[1,1]; a=[F(0),F(1,3)]; b=[F(0),F(1)]
for j in range(2,196):
    m.append(((2*j+1)*m[-1]+3*(j-1)*m[-2])//(j+2))
    a.append(((2*j+1)*a[-1]+3*(j-1)*a[-2]
              +2*m[-2]+3*m[-3]-m[-1])/(j+2))
    b.append(((2*j+1)*b[-1]+3*(j-1)*b[-2])/(j+2))
def res(t):
    t=F(t)
    return t.numerator*pow(t.denominator,-1,64)%64
assert [res(m[-1]),res(16*a[-1]),res(16*b[-1])]==[52,7,11]
motz = [1, 1]
for j in range(2, 2**14 + 196):
    motz.append(((2*j + 1)*motz[-1] + 3*(j - 1)*motz[-2])//(j + 2))
def vdiff(n):
    d = motz[2**(n+1) + 195] - motz[2**n + 195]
    return (d & -d).bit_length() - 1
assert [vdiff(11), vdiff(13)] == [12, 16]
print('all identities verified')
\end{verbatim}}
\section{A linear range of simultaneous orders}
\label{sec:linear}
The following statement quantifies the dependence of the initial exceptions on the order. Write $\Delta a_n=a_{n+1}-a_n$. For a finite measure $\mu$ on a bounded real interval, let $a_n=\int u^n\,d\mu(u)$. We use a local condition only at the endpoint of largest modulus.

\begin{theorem}[Endpoint moments]\label{thm:endpoint}

Suppose the largest endpoint is $\lambda>0$, the remaining support is contained in $[-q\lambda,\lambda]$ with $0\le q<1$, and in a left neighborhood of $\lambda$ the pushforward by $t=-\log(u/\lambda)$ has density $t^{\alpha-1}h(t)$, where $h$ is smooth up to $0$, $h(0)>0$, and $\alpha>0$ is nonintegral. Then there are $c>0$ and $n_0$ such that
\[
 (-1)^r\Delta^r\log a_n>0\qquad(n\ge n_0,\quad 2\le r\le cn).
\]
The statement still holds for a finite signed measure satisfying the same local hypothesis. In particular, its moments are eventually positive.\end{theorem}

\begin{proof}
Choose a smooth cutoff $\chi(t)$ equal to $1$ near $0$, nonnegative and compactly supported in the local coordinate interval, where $h(t)>0$. Set $f(t)=\chi(t)t^{\alpha-1}h(t)$ for $t>0$, and
$F(x)=\int_0^\infty e^{-xt}f(t)\,dt$. Then
$\lambda^{-n}a_n=F(n)+O(e^{-bn})$ for some $b>0$, since outside the retained endpoint neighborhood the absolute value of $u/\lambda$ is uniformly smaller than $1$. The substitution $y=nt$ and dominated convergence give $F(n)\sim A n^{-\alpha}$, where $A=h(0)\Gamma(\alpha)>0$. Hence for some $\kappa>0$,
\[
 \log a_n=n\log\lambda+\log F(n)+O(e^{-\kappa n}).
\]

Write $m=\lceil\alpha\rceil$, $\beta=m-\alpha\in(0,1)$, and
$R(t)=f(t)-h(0)t^{\alpha-1}e^{-t}$. Near zero $R(t)$ has the form $t^\alpha$ times a smooth function. Thus $R^{(j)}(0+)=0$ for $0\le j<m$, and $(D+1)^m R(t)=O(t^{\alpha-m})$ near zero, with $D=d/dt$. It decays exponentially times a power at infinity. Let
\[
 b_\beta(t)=\frac{t^{\beta-1}e^{-t}}{\Gamma(\beta)},\qquad
 g=A^{-1}b_\beta * (D+1)^mR,
\]
where $*$ is convolution on $[0,\infty)$. This $g$ is bounded: near zero this follows from $\int_0^t(t-s)^{\beta-1}s^{-\beta}\,ds=\mathrm B(\beta,1-\beta)$, and away from zero one splits the integral at $t/2$ and uses local integrability and exponential decay. For real $x>0$, integration by parts $m$ times (all endpoint terms at zero vanish) and the Laplace convolution formula give
\[
 F(x)=A(x+1)^{-\alpha}\bigl(1+\mathcal L g(x)\bigr),
 \qquad \mathcal Lg(x)=\int_0^\infty e^{-xt}g(t)\,dt .
\]

Fix $B\ge\max(1,\|g\|_\infty)$. The convolution bound $|g^{*j}(t)|\le B^j t^{j-1}/(j-1)!$ implies that
$k(t)=\sum_{j\ge1}(-1)^{j+1}g^{*j}(t)/j$ satisfies $|k(t)|\le B e^{Bt}$. For $x>B$, absolute convergence and the real logarithm series imply
$\log(1+\mathcal Lg(x))=\mathcal Lk(x)$. Differentiating under the integral gives
\[
 |(\mathcal Lk)^{(r)}(x)|\le B\frac{r!}{(x-B)^{r+1}}.
\]
Relative to $\alpha(r-1)!/(x+1)^r$, the magnitude of this remainder is bounded by
\[
 \frac{Br}{\alpha(x-B)}
 \exp\!\left(\frac{r(B+1)}{x-B}\right).
\]
For $x\ge 2B$ and $1\le r\le c_0 x$, this is less than $1/2$ if $c_0>0$ is chosen sufficiently small. As $\log F=\log A-\alpha\log(x+1)+\mathcal Lk$, integration of $(-1)^r(\log F)^{(r)}$ over the translated unit $r$-cube yields
\[
 (-1)^r \Delta^r\log F(n)\ \ge\ \frac{\alpha(r-1)!}{2(n+r+1)^r}
 \qquad(1\le r\le c_0 n,\ n\ge 2B).
\]
The difference between the original logarithmic difference and this cutoff expression is $O(2^r e^{-\kappa n})$ for $r\ge 2$. To compare the two bounds, write $c\le\min(1,c_0)$. Since $(r-1)!\ge (r/e)^r/r$ and $n+r+1\le 3n$, the lower bound exceeds $\alpha(2n)^{-1}\exp[-cn\log(3e/c)]$ uniformly for $1\le r\le cn$; here $t\mapsto t\log(3e/t)$ is increasing on $(0,1]$. Choosing $c$ still smaller so that $c\log(6e/c)<\kappa$ completes the proof.\end{proof}
\section{The five sequences of Chen, Guo and Wang}\label{sec:five}
We verify the hypotheses of the endpoint theorem one sequence at a time.
All the endpoints below refer to the variable being raised to the $n$-th power.
\begin{center}
\small\begin{tabular}{@{}lrrl@{}}
\toprule
Sequence & Upper endpoint $\lambda$ & Exponent $\alpha$ & Moment model\\
\midrule
Motzkin & $3$ & $3/2$ & translated semicircle\\
Fine & $4$ & $3/2$ & continuous part and a negative atom\\
Central Delannoy & $3+2\sqrt2$ & $1/2$ & arcsine law\\
Tree-like polyhexes & $5$ & $3/2$ & translated semicircle with an index shift\\
Domb & $16$ & $3/2$ & three-dimensional angular integral\\
\bottomrule
\end{tabular}
\end{center}

\paragraph{Motzkin and tree-like polyhex numbers.}
Let $m_j(a)$ be the $j$-th moment of the probability density
$(2\pi)^{-1}\sqrt{4-(u-a)^2}$ on $[a-2,a+2]$. Substitution in the beta integral shows that
the even moments about $a$ are $C_j$ at order $2j$, and that odd central moments vanish.
Consequently $m_j(1)=M_j$. Writing $H_a(z)=\sum_{j\ge0}m_j(a)z^j$ gives
\[
 H_a(z)=\frac{1-az-\sqrt{(1-az)^2-4z^2}}{2z^2}.
\]
The original polyhex sequence has $t_0=t_1=1$ and
$(n+1)t_n=3(2n-1)t_{n-1}-5(n-2)t_{n-2}$ for $n\ge2$
\cite[Sec.~4]{cgw}. To check the index, put $h_j=m_j(3)$. Differentiating
the displayed generating function or comparing coefficients gives
$(j+2)h_j=3(2j+1)h_{j-1}-5(j-1)h_{j-2}$ for $j\ge2$, with
$h_0=1,h_1=3$. Thus $t_n=m_{n-1}(3)$ for $n\ge1$. One may equally use
the measure $u^{-1}\,d\mu_3(u)$ on $[1,5]$ to represent $t_n$ directly
on that tail. The right endpoint in either semicircle case is a square-root
zero. The change $u=\lambda e^{-t}$ therefore gives $t^{1/2}$ times a
smooth positive factor at zero.

\paragraph{Central Delannoy numbers.}
Their generating function is $(1-6z+z^2)^{-1/2}$ \cite{cgw}.
For $a=3-2\sqrt2$ and $b=3+2\sqrt2$, the arcsine probability
density $[\pi\sqrt{(u-a)(b-u)}]^{-1}$ on $(a,b)$ has precisely this
moment series. Indeed its integral against $(1-zu)^{-1}$ is
$[(1-az)(1-bz)]^{-1/2}$, by the cosine substitution. Its density has a
right endpoint singularity of order $-1/2$.

\paragraph{Fine numbers.}
Let $F_n$ have $F_0=1$, $F_1=0$ and $2F_n+F_{n-1}=C_n$ for $n\ge1$
\cite{fine}. Let $\nu$ be the semicircle probability law centered at $2$,
with density $\sqrt{u(4-u)}/(2\pi)$ on $(0,4)$. Its moments are
$C_{n+1}$, by a beta integral. Set
\[
 d\mu(u)=\tfrac34\,d\delta_{-1/2}(u)+(2u+1)^{-1}d\nu(u).
\]
This measure has mass one: the transform
$\int(w-u)^{-1}d\nu(u)=(w-2-\sqrt{(w-2)^2-4})/2$,
with the square root asymptotic to $w-2$, is $-1/2$ at $w=-1/2$.
The continuous mass is thus $1/4$. Multiplying the measure by $2u+1$
removes the atom and gives $\nu$, proving that its moments satisfy the
stated Fine recurrence. Its upper edge is again a square-root zero.
The single zero $F_1$ causes no restriction on logarithms taken on a
positive tail.

\paragraph{Domb numbers.}
In the usual binomial sum for the Domb numbers, Vandermonde's identity
in each of two pairs of variables gives
\[
 D_n=\sum_{k=0}^n\binom nk^2\binom{2k}{k}\binom{2n-2k}{n-k}
 =\!\sum_{k_1+\cdots+k_4=n}\!\binom{n}{k_1,k_2,k_3,k_4}^{\!2}.
\]
The angular model also appears as the even moments of a four-step planar
random walk \cite[Example~3.17]{walk}. Put
$V(\theta)=|1+e^{i\theta_1}+e^{i\theta_2}+e^{i\theta_3}|^2$
on the three-torus with its probability Haar measure.
Fourier orthogonality gives $D_n=\int V^n$.
Equality in the triangle inequality forces the unique maximum $V=16$
to occur at zero. There
\[
 16-V(\theta)=4\sum_j\theta_j^2-\Big(\sum_j\theta_j\Big)^2
                +O(|\theta|^4).
\]
The quadratic form is positive definite. The analytic Morse lemma gives
coordinates $y$ in which $t=\log(16/V)=|y|^2$.
The Jacobian in these coordinates is analytic and positive near zero,
after choosing orientation. In polar coordinates its spherical average
at radius $\sqrt t$ is analytic in $t$ (odd homogeneous terms integrate
to zero) and has positive constant term. Thus the pushforward density
is $t^{1/2}h(t)$ with $h(0)>0$. On the compact complement of any such
neighborhood, $V/16$ is bounded away from one.

\begin{corollary}\label{cor:five}
There exist $c>0$ and $N_0$ common to the five sequences in the table
such that $S_{n,r}(a)>0$ whenever $n\ge N_0$ and $2\le r\le cn$.
For each fixed $r\ge2$ one also has
$S_{n,r}(a)\sim\alpha(r-1)!\,n^{-r}$.
\end{corollary}
\begin{proof}
Apply the endpoint theorem and take the minimum of five slopes and
the maximum of five thresholds. For fixed $r$ the logarithmic derivative
remainder in its proof is $O_r(n^{-r-1})$ throughout the cube of
integration. The exponentially small discarded contribution is smaller
still.
\end{proof}

For comparison, each density in the table has an analytic factor at the
edge, hence expansion to any fixed order after Laplace integration:
expand that factor at $0$ with Taylor remainder and integrate the monomials.
This gives $a_n=K\lambda^n n^{-\alpha}(1+\sum_{j=1}^J c_j n^{-j}
+O_J(n^{-J-1}))$ for every fixed $J$. In particular
$a_{n-1}a_{n+1}/a_n^2=1+\alpha n^{-2}+O(n^{-3})$ and its expansion
continues to arbitrary algebraic order. These are precisely the data
required by Li's fixed-order criterion \cite[Thm.~2.2]{li}. Thus the
fixed-order conjecture follows independently by a pre-existing
criterion once its hypotheses are checked. The uniform estimate in
Corollary~\ref{cor:five} also controls orders that increase with $n$.

\section{The critical order for the Motzkin numbers}\label{sec:motzkin}
Let $S_{n,r}=(-1)^r\Delta^r\log M_n$. The Motzkin numbers admit the semicircle moment formula
\[
 M_n=\frac1{2\pi}\int_{-1}^3 u^n\sqrt{(3-u)(u+1)}\,du
    =P_n+(-1)^nQ_n,
\]
where $P_n$ is the integral over $[0,3]$, and $Q_n$ is the same positive integral over $[-1,0]$ with $u^n$ replaced by $|u|^n$. The moment formula also follows directly by expanding $(1+x)^n$ against the centered semicircle law on $[-2,2]$, whose even moments are Catalan numbers and odd moments vanish.

\begin{theorem}[Signs in a Motzkin row]\label{thm:transition}
For every $0<c<1/2$, $S_{n,r}>0$ for all sufficiently large $n$ and all integers $2\le r\le cn$. For every $c>1/2$, all sufficiently large $n$ satisfy $\operatorname{sign}S_{n,r}=(-1)^n$ for every $r\ge cn$. The same parity conclusion holds for $r=n/2+O(1)$. On all sufficiently large even $n$, $S_{n,r}>0$ for every $r\ge2$. On odd $n$, let $r_*(n)$ be the first integer $r\ge2$ with $S_{n,r}\le0$. For all sufficiently large odd $n$, every $r>r_*(n)$ satisfies $S_{n,r}<0$, and
\[
 r_*(n)=\frac n2-\frac{\log n}{2\log4}+O(1).
\]

\end{theorem}
\begin{proof}
Substituting $u=3e^{-t}$ gives $3^{-n}P_n=F(n)=\mathcal L f(n)$ where \[ f(t)=\frac{3e^{-t}}{2\pi}\sqrt{3(1-e^{-t})(1+3e^{-t})}. \] Set $\alpha=3/2$, $a=13/8$, $C=3\sqrt3/\pi$ and
\[
 f(t)=Ct^{1/2}e^{-at}\phi(t),\qquad
 \phi(t)=e^{5t/8}\sqrt{X(t)Y(t)},\quad
 X(t)=\frac{1-e^{-t}}t,\quad Y(t)=\frac{1+3e^{-t}}4.
\]
Use the construction of the endpoint theorem with reference density $Ct^{1/2}e^{-at}$ and Laplace shift $a$ instead of $1$. This applies on the whole positive half-line because $f$ and its first two derivatives decay exponentially at infinity. Writing $I^{1/2}v(t)=\Gamma(1/2)^{-1}\int_0^t(t-s)^{-1/2}v(s)\,ds$, the two integral kernels in that construction become
\[
 g(t)=e^{-at}G(t),\quad G=\frac{I^{1/2} D^2[t^{1/2}(\phi-1)]}{\Gamma(3/2)},\quad
 k(t)=e^{-at}K(t),\quad K(t)=\sum_{j\ge1}\frac{(-1)^{j+1}}jG^{*j}(t).
\]
For $n$ large and $r\ge2$, the integral identity for the logarithm and the convolution series give
\[
 (-1)^r\Delta^r\log P_n
 =\int_0^\infty e^{-nt}(1-e^{-t})^r H(t)\,dt,
 \qquad H(t)=\frac{\alpha e^{-at}}{t}+k(t).
\]
Here $|k(t)|\le B e^{Bt}$ for a fixed $B$.

We require positivity of $H$ on a specific interval. Let $J=\log\phi$. The functions $X,Y$ are Laplace transforms of the uniform law on $[0,1]$ and the Bernoulli law with success probability $3/4$, respectively. Thus $J(0)=J'(0)=0$, and $J''$ equals half the sum of the variances under the two exponential tilts. Every random variable supported in $[0,1]$ has variance at most $1/4$, since $\operatorname{Var}(Z)\le\mathbb E[(Z-1/2)^2]\le1/4$. Hence $0\le J''\le1/4$ and on $0\le t\le1/2$,
\[
 |\phi''(t)|\le e^{1/32}(1/4+1/64)<1/3,\qquad
 |\phi'(t)|\le t/3,\qquad |\phi(t)-1|\le t^2/6.
\]
The product rule bounds $|D^2[t^{1/2}(\phi-1)]|$ by $(17/24)t^{1/2}$. As $I^{1/2}t^{1/2}=\Gamma(3/2)t$, we obtain $|G(t)|\le17t/24<1/2$ on this interval. Convolution powers there satisfy $|G^{*j}(t)|\le 2^{-j}t^{j-1}/(j-1)!$, so $|K(t)|\le e^{1/4}/2<1$. It follows that
\[
 H(t)=e^{-at}(\alpha/t+K(t))>0\qquad (0<t\le1/2).
\]

For $c=r/n$ let $p_c(t)=-t+c\log(1-e^{-t})$. Its unique maximizer is $t_c=\log(1+c)$, its value there is $L(c)=c\log c-(1+c)\log(1+c)$, and $p_c''(t_c)=-(1+c)/c$. If $c$ ranges over a compact interval in $(0,e^{1/2}-1)$, the positive saddle yields, uniformly,
\[
 (-1)^r\Delta^r\log P_n=n^{-1/2}e^{nL(c)}\bigl(D(c)+o(1)\bigr),\qquad
 D(c)=H(t_c)\sqrt{2\pi c/(1+c)}>0.
\]
For completeness, restrict to a fixed radius around each $t_c$ contained in $(0,1/2)$. Taylor expansion on the $n^{-2/5}$ neighborhood gives the Gaussian integral uniformly because $n(n^{-2/5})^3\to0$. From $n^{-2/5}$ to the fixed radius, strict negative curvature gives a factor $e^{-b n^{1/5}}$ relative to $e^{nL(c)}$. Off the fixed neighborhoods, the maximum drops by a uniform $\eta>0$. The singularity at zero and the exponential bound at infinity are absorbed by the integrable factor $e^{m p_c(t)}|H(t)|$ with $m>B+1$ fixed. These arguments also give the corresponding uniform upper bound $O(n^{-1/2}e^{nL(c)})$ on every compact subinterval of $(0,\infty)$, regardless of the sign of $H(t_c)$.

The endpoint estimates at $3$ and $-1$ give
\[
 d_n=Q_n/P_n=d\,3^{-n}(1+O(1/n)),\qquad d=(3\sqrt3)^{-1}.
\]
Indeed, the densities in $u=3e^{-t}$ and $|u|=e^{-t}$ on the respective intervals are $t^{1/2}$ times smooth factors with constant terms $3\sqrt3/\pi$ and $1/\pi$. Rescaling $t$ by $n^{-1}$ gives the stated expansions with a relative error $O(1/n)$. Since $\log M_n=\log P_n+(-1)^nd_n+O(d_n^2)$, the binomial difference formula now gives uniformly in integer $r\ge2$
\[
 S_{n,r}=(-1)^r\Delta^r\log P_n
    +(-1)^n d\,3^{-n}(4/3)^r(1+O(1/n))
    +O(9^{-n}(10/9)^r).
\]
The error is exponentially small compared with $3^{-n}(4/3)^r$ for all such $r$. Set $A(c)=-\log3+c\log(4/3)$. Then $L(1/2)=A(1/2)$, and $(A-L)'(c)=\log(4(1+c)/(3c))>0$. The saddle estimate establishes positivity uniformly on compact subintervals of $(0,1/2)$ and shows parity dominance for $r/n\to c>1/2$. The endpoint theorem covers all small $r/n$ down to $r=2$, proving the full subcritical assertion. The extra factor $n^{-1/2}$ in the saddle contribution proves parity dominance for $r=n/2+O(1)$. For $r\ge4n$, $3^{-n}(4/3)^r\ge(256/243)^n$, whereas the absolute value of the saddle contribution is bounded independently of $n,r$, by inserting $(1-e^{-t})^2$ in place of $(1-e^{-t})^r$ and integrating $e^{-nt}(1-e^{-t})^2|H(t)|$. Compact uniform estimates on $r/n\in[1/2+\delta,4]$ handle the remaining supercritical orders. In a neighborhood of $1/2$, the saddle contribution is positive and so is the oscillating contribution for even $n$, which proves positivity for all orders in every sufficiently large even row.

To locate the first exceptional order, put $D_*=D(1/2)>0$. Near $c=1/2$, the ratio of the positive saddle to the magnitude of the oscillating term is
\[
 \frac{D(c)}d n^{-1/2}\exp\{n[L(c)-A(c)]\}(1+o(1)).
\]
The function $L-A$ decreases strictly and has derivative $-\log4$ at $1/2$. For $r=n/2+O(\log n)$ the ratio equals
\[
 (1+o(1))4^{T_n-r},\qquad
 T_n=\frac n2-\frac{\log n}{2\log4}+\frac{\log(D_*/d)}{\log4}.
\]
The estimate is uniform on each such logarithmic window, by Taylor's formula with quadratic error $O((\log n)^2/n)$. For odd $n$, all orders $2\le r<T_n-\eta$ are good and the first integer $r>T_n+\eta$ is bad, for any fixed $\eta>0$ and all sufficiently large $n$. In fact the compact subcritical range and the endpoint theorem handle $r/n$ separated from $1/2$. Inside a fixed neighborhood of $1/2$ but below the logarithmic window, monotonicity of $L-A$ with derivative bounded away from zero makes the ratio diverge uniformly if the window is chosen sufficiently wide. Inside the window the displayed ratio gives both claims. The same monotonicity gives ratios tending uniformly to zero above the logarithmic window, within the fixed neighborhood. In the window the oscillating contribution including its relative error has sign $(-1)^n$ and its magnitude divided into the positive saddle is $(1+o(1))4^{T_n-r}$. Hence consecutive ratios decrease by a factor tending uniformly to $1/4$. For odd $n$ there is a single change from positive to nonpositive, with all subsequent entries strictly negative. This proves the assertions about $r_*(n)$.\end{proof}
\section{A single change of log-curvature on every vertical Delannoy line}\label{sec:vertical}

Write $D(i,j)$ for the number of walks from $(0,0)$ to $(i,j)$ with steps
$(1,0),(0,1),(1,1)$, and put $d(i,j)=D(i-j,j)$. Wang, Zheng and Chen
asked whether $d(h+2k,k)$ is eventually log-convex for each fixed $h\ge0$.
Their Property~(B) proposes a single change from log-concavity to log-convexity
on all interior rays \cite[Sec.~4]{wzc}. We establish that sign pattern on
every vertical line, for the original sequence $u_k=d(h+2k,k)=D(h+k,k)$.

\begin{theorem}\label{thm:vertical}
Let $T_k=u_{k-1}u_{k+1}-u_k^2$ for $k\ge1$. For $h=0$, every $T_k$ is
positive. For $h=1$, $T_1=0$ and every later $T_k$ is positive. For every
integer $h\ge2$, there exists $j\ge2$ such that $T_k<0$ for $1\le k<j$,
$T_j\ge0$, and $T_k>0$ for $k>j$.
\end{theorem}

\begin{proof}
We use $u_k=P_k^{(h,0)}(3)$, the unit-weight specialization of Hetyei's
Jacobi identity \cite[Prop.~2.9]{hetyei}. It follows directly from
$D(h+k,k)=\sum_i2^i\binom{k}{i}\binom{h+k}{i}$ and the binomial formula
for $P_k^{(h,0)}(3)$. For $h=0$, $u_k$ is the $k$th moment of
$3+2\sqrt{2}\cos\theta$ under uniform measure on $[0,\pi]$, by the
Legendre generating function. Strict Cauchy--Schwarz gives $T_k>0$.

For $h\ge1$ define $R_k=u_k/u_{k-1}$ and $E_k=R_{k+1}-R_k$ for
$k\ge1$. The signs of $E_k$ and $T_k$ coincide. The three-term Jacobi
recurrence at $x=3$ gives $R_{k+1}=A_k-B_k/R_k$ for $k\ge1$, where,
with $s=2k+h$,
\[
 A_k=\frac{(s+1)[3s(s+2)+h^2]}{2(k+1)(k+h+1)s},\qquad
 B_k=\frac{k(k+h)(s+2)}{(k+1)(k+h+1)s}.
\]
These formulas also give $A_0=R_1=2h+3$ and $B_0=0$. Set
$H_k=(A_{k+1}-A_k)/(B_{k+1}-B_k)$ for $k\ge0$. Subtraction gives,
for every $k\ge1$,
\begin{equation}\label{eq:vertical-propagate}
 E_{k+1}=(B_{k+1}-B_k)(H_k-1/R_{k+1})
          +\frac{B_k E_k}{R_kR_{k+1}}.
\end{equation}
To determine the monotonicity of $H_k$, put $z=2k+h+2$ and
$V(z)=z^2(z^2-4)+h^2(4z^2-h^2)>0$. Direct subtraction in the rational
formulas above gives
\begin{align*}
 B_{k+1}-B_k&=\frac{4V(z)}{z(z^2-h^2)(z-2)((z+2)^2-h^2)},\\
 H_k&=3-\frac{4h^2(2z^3+z^2-2z-h^2)}{V(z)},\\
 H_{k+1}-H_k&=\frac{16h^2 z(z+1)((z+2)^2-h^2)(z^2+2z-3h^2)}
                      {V(z)V(z+2)}.
\end{align*}
Thus $B_{k+1}>B_k$, and the sign of $H_{k+1}-H_k$ is that of
$z^2+2z-3h^2$, strictly increasing with $k$. The initial values are
\[
 H_0=-\frac{2h^3+7h^2+h-12}{(h+1)(h+4)},\qquad
 E_1=-\frac{2(h-1)(h+2)}{2h+3}.
\]
For $h\ge2$, $H_0<0$. Consequently, as soon as $H_i>0$, all increments
$H_{k+1}-H_k$ with $k\ge i$ are strictly positive. Otherwise every increment
up to index $i$ would have been nonpositive, forcing $H_i\le H_0$.

For completeness, the possibility of $E_k<0$ for every $k$ can already be
excluded by a leading-term estimate. Set $\Lambda=3+2\sqrt2$. The exterior
Jacobi asymptotic gives $u_k\sim C_h\Lambda^k k^{-1/2}$ with $C_h>0$
\cite[Thm.~8.21.9]{szego}. Thus $R_k\to\Lambda$ and $u_k/\Lambda^k\to0$.
If all $E_k<0$, then $R_k>\Lambda$ for every $k$, and
$u_k/\Lambda^k=\prod_{i=1}^k(R_i/\Lambda)\ge1$, a contradiction.
Now suppose $h\ge2$ and take the least $j$ with $E_j\ge0$. Since $E_1<0$,
we have $j\ge2$. Equation~\eqref{eq:vertical-propagate} at $k=j-1$ and
$E_{j-1}<0$ give $H_{j-1}>1/R_j>0$. Hence
$H_j>H_{j-1}>1/R_j\ge 1/R_{j+1}$, so $E_{j+1}>0$. The same inequality
propagates to every later index. If $h=1$, then $H_0=1/5=1/R_1$, $E_1=0$,
and $H_k$ is strictly increasing from $k=0$. Equation~\eqref{eq:vertical-propagate}
gives $E_2>0$ and proves strict positivity at all later indices by induction.
\end{proof}

The fixed-$h$ eventual assertion also follows by taking three terms in the
classical exterior Jacobi expansion. The recurrence estimate
\eqref{eq:vertical-propagate} controls every preceding index as well. The
one-change assertion for rays of arbitrary interior slope in \cite[Sec.~4]{wzc}
extends beyond Theorem~\ref{thm:vertical}.
\section{Sources, scope and priority}\label{sec:audit}
Lengyel's Conjectures~5.1--5.3 are printed together in
\cite[pp.~132--133]{lengyel13}. Sections~\ref{sec:lengyel} and
\ref{sec:paired} establish 5.3 and 5.2 with the original quantifiers.
The counterexample in Section~\ref{sec:counter} contradicts a universal
upper bound stated in 5.1. Its lower bound has the effective form
$n-\lfloor\log_2(b+2)\rfloor$ for all odd multipliers.
The existence of a finite upper correction depending on each pair $(a,b)$
remains unresolved here.

Lengyel's subsequent paper \cite{lengyel14} determines orders of the
unshifted prime-power differences. Its binary calculation is a substantive
input here. Krattenthaler and M\"uller \cite{km18} determine individual
Motzkin numbers modulo eight. Rowland--Yassawi and Straub give general ways to compute fixed-modulus
automata \cite{ry,straub}. The same modulo-eight classification was
obtained independently by Wang and Xin, as recorded in the note following
Theorem~17 of \cite{km18}. The passage from these residues to valuations
of differences modulo powers growing with $n$, uniformly over the
family of shifts in Conjecture~5.2 in the sense $\forall q\,\exists n_0(q)$,
requires the additional coefficient calculation above.

The methodological starting points in VMRC \cite{corpus} are identifiable.
Manuscript [006] extracts a torus moment and a definite Hessian from a
constant term. Manuscript [011] separates the contribution near a dominant
endpoint from an exponentially small complementary integral. Manuscript
[070] tracks valuations after normalization in a Frobenius-lifting
problem. Here the analytic application needs a bound after taking
arbitrarily high logarithmic differences; the arithmetic application needs
the residues of the derivatives of the shift recurrence. The proofs above
give both estimates directly. They invoke no unverified theorem from the corpus.

For the analytic problem Zhu \cite{zhu} treats several low orders and
settles all orders for derangements. Li's criterion \cite{li} gives a
prior route to each fixed order for the five sequences once their endpoint
expansions are supplied. Wang and Yang \cite{wy} already study
growing-order differences for the partition function and for coefficients
of the Riemann Xi function. The specific claims to compare for the present
analysis are the five-sequence linear range and the Motzkin first-sign-change
formula. Their proofs control an exponentially smaller part of a moment
that becomes decisive after taking sufficiently many differences.

\begin{center}\small
\begin{tabularx}{\linewidth}{@{}p{.13\linewidth}X@{}}
\toprule
Criterion & Finding for Lengyel's Conjectures 5.1--5.3 \\
\midrule
G1, G2 & The named 2013 source specifies $M_n$, the offsets and all
quantifiers. Sections~\ref{sec:lengyel}, \ref{sec:paired} and \ref{sec:counter} use those definitions.\\
G3 & The author records a conditional route for Conjecture~5.2 and the
known cases of Conjecture~5.3. The later publications checked do not
supply the full paired law, the three-shift conclusion or the offset-195 counterexample. This describes the literature
examined through 15 November 2026.\\
G4 & The arithmetic proofs include the small exponents and both parities;
the sharp cutoff for $q\ge4$ is $q+1$. Exact polynomial identities and
the Taylor tail bound prove the counterexample for every odd $n\ge13$.\\
G5 & The dyadic derivative calculation detects the cancellation after
doubling the index for every power-of-two shift and at the exceptional
offset $195$. A calculation with a fixed
small modulus for the original terms cannot on its own certify these orders.\\
G6 & The search covered the original and follow-up papers, the two modulo-eight
classifications, author/title and formula searches, bibliographic matches
and the supplied corpus. No earlier proof or disproof of the three complete claims was located. Historical
priority and independent external verification remain unconfirmed.\\
\bottomrule
\end{tabularx}
\end{center}
\begin{thebibliography}{99}\small
\bibitem{lengyel13} T.~Lengyel, On divisibility properties of some differences of Motzkin numbers, \emph{Ann. Math. Inform.} \textbf{41} (2013), 121--136. \href{https://sites.oxy.edu/lengyel/papers/AMI_41_from121to136.pdf}{Author's copy}.
\bibitem{lengyel14} T.~Lengyel, Exact $p$-adic orders for differences of Motzkin numbers, \emph{Int. J. Number Theory} \textbf{10} (2014), 653--667. \href{https://doi.org/10.1142/S1793042113501157}{doi:10.1142/S1793042113501157}.
\bibitem{km18} C.~Krattenthaler and T.~W.~M\"uller, Motzkin numbers and related sequences modulo powers of 2, \emph{European J. Combin.} \textbf{73} (2018), 114--137. \href{https://arxiv.org/abs/1608.05657}{arXiv:1608.05657}.
\bibitem{ry} E.~Rowland and R.~Yassawi, Automatic congruences for diagonals of rational functions, \emph{J. Th\'eor. Nombres Bordeaux} \textbf{27} (2015), 245--288. \href{https://doi.org/10.5802/jtnb.901}{doi:10.5802/jtnb.901}.
\bibitem{straub} A.~Straub, On congruence schemes for constant terms and their applications, \emph{Res. Number Theory} \textbf{8} (2022), article 42. \href{https://arxiv.org/abs/2205.09902}{arXiv:2205.09902}.
\bibitem{cgw} W.~Y.~C.~Chen, J.~J.~F.~Guo and L.~X.~W.~Wang, Infinitely log-monotonic combinatorial sequences, \emph{Adv. Appl. Math.} \textbf{52} (2014), 99--120. \href{https://arxiv.org/abs/1304.5160}{arXiv:1304.5160}.
\bibitem{li} G.-J.~Li, $\ell$-Log-momotonic and Laguerre inequality of P-recursive sequences, 2022, Theorem~2.2. \href{https://arxiv.org/abs/2206.13922}{arXiv:2206.13922}.
\bibitem{zhu} B.-X.~Zhu, Higher order log-monotonicity of combinatorial sequences, 2013. \href{https://arxiv.org/abs/1309.6025}{arXiv:1309.6025}.
\bibitem{fine} E.~Deutsch and L.~Shapiro, A survey of the Fine numbers, \emph{Discrete Math.} \textbf{241} (2001), 241--265. \href{https://www.sciencedirect.com/science/article/pii/S0012365X01001212}{Publisher's article}.
\bibitem{walk} J.~M.~Borwein, A.~Straub and C.~Vignat, Densities of short uniform random walks in higher dimensions, \emph{J. Math. Anal. Appl.} \textbf{437} (2016), 668--707. \href{https://arxiv.org/abs/1508.04729}{arXiv:1508.04729}.
\bibitem{wy} L.~X.~W.~Wang and N.~N.~Y.~Yang, Laguerre inequalities and complete monotonicity for the Riemann Xi-function and the partition function, \emph{Trans. Amer. Math. Soc.} \textbf{377} (2024), 4703--4725. \href{https://doi.org/10.1090/tran/9081}{doi:10.1090/tran/9081}.
\bibitem{wzc} Y.~Wang, S.-N.~Zheng and X.~Chen, Analytic aspects of Delannoy numbers, \emph{Discrete Math.} \textbf{342} (2019), 2270--2277. \href{https://doi.org/10.1016/j.disc.2019.04.003}{doi:10.1016/j.disc.2019.04.003}.
\bibitem{hetyei} G.~Hetyei, Shifted Jacobi polynomials and Delannoy numbers, Proposition~2.9. \href{https://arxiv.org/abs/0909.5512}{arXiv:0909.5512}.
\bibitem{szego} G.~Szeg\H{o}, \emph{Orthogonal Polynomials}, 4th ed., AMS Colloquium Publications, vol.~23 (1975), Theorem~8.21.9.
\bibitem{corpus} QEEDDD, \emph{Verifiable Mathematical Reasoning Corpus}, manuscripts [006], [011] and [070], repository snapshot
\href{https://github.com/QEEDDD/Verifiable-Mathematical-Reasoning-Corpus/tree/5887cdea4898f0a3aa5bf7a8e0829697bd2b0f4b}{5887cdea4898f0a3aa5bf7a8e0829697bd2b0f4b}.
\end{thebibliography}
\end{document}
